\section{Related Work}
\label{sec:RelatedWork}

%% Iverson's languages
Rank polymorphism originally appeared in APL \cite{APL},
which Iverson designed as a form of mathematical notation,
with the APL interpreter serving to eliminate the semantic ambiguity
found in conventional notation \cite{IversonTuringLecture}.
At first, APL only lifted \emph{scalar} functions
to operate on aggregate data via pointwise application,
either on a scalar argument and an aggregate argument
or on two aggregates of identical shape.
Subsequent development introduced the notion of function rank,
the number of dimensions a function expects its argument to have.
This generalized the scalar function lifting,
\eg, allowing a vector-mean function to produce a vector of results
when given a matrix argument,
and it introduced the ``frame of cells'' view of aggregate arguments
where pointwise lifting generalizes to cellwise lifting.
The next generalization step was to loosen the rule
on frame shape compatibility.
In J \cite{J}, which Iverson created as a successor to APL,
a function can be applied to two arguments of differing frames
as long as one frame---%
viewed as a sequence of dimensions---%
is a prefix of the other.
This was a conscious design decision on Iverson's part:
prefix agreement was chosen over suffix agreement because
it fit better with APL's emphasis on operating along arrays' major axes \cite{RankAndUniformity}.

%% other array languages
FISh also made implicit aggregate lifting part of the semantics of function application,
and its static semantics resolves the shapes of all arrays \cite{FISh}.
A conventional type judgment describes
the elements of arrays computed by a FISh program,
and a second judgment ascribes a shape to each array.
In FISh's metatheory,
the shape of a function is a function on shapes.
Thus a function application's shape can be calculated statically
by applying the shape function to the arguments' shapes.

As elegant as this model is in a first-order language,
it is incompatible with first-class functions.
When functions appear in arrays which are themselves applied to arguments,
the shape must describe
the layout of of that collection of functions,
not just the functions' own behavior.
In Remora, a function which checks
whether a point in $\mathbb{R}^3$ is near the origin
might have the type
{\tt (-> ((Arr Float (Shp 3))) (Arr Bool (Shp)))}.
FISh considers this function's shape to be
the function on a singleton domain
(containing only $[3]$)
which returns the empty vector.
However, Remora expressions produce array data,
even in function position.
A function array {\tt near-origin?} containing only that function
has its own ``first-order'' shape {\tt (Shp)},
independent of any function on shapes summarizing its behavior.

In resolving all shapes statically,
FISh is also too restrictive to permit shapes determined from run-time data.
Functions like {\tt iota} and {\tt filter} cannot exist,
nor can the ragged data which would result from lifting them.
By characterizing functions with restricted dependent types,
Remora escapes both of these limitations.

%%% \cite{Transfinite} arrays (Sinkarovs & Scholz)

%%% \cite{SISAL}

ZPL is a data-parallel language
which was designed to live within a larger language
with a more general parallelism mode \cite{ZPLArraySublang}.
Its programming model is based on
an explicit map operation
over programmer-specified index space within an array.
Several built-in operators modify an index space,
such as shifting a section along some dimension,
adding a new dimension by broadcasting,
or slicing out a particular sub-array.
The set of built in operators is constructed
to make communication cost implications clear to the programmer \cite{ZPLPerformance}.

NESL uses a nested-vector data model,
rather than the rank-polymorphic regular-array model \cite{NESL}.
Programs map operations over vectors using a comprehension notation.
Since a vector's elements may themselves be vectors
with widely varying lengths,
NESL's main performance trick is
turning irregular nested vectors into a flat internal representation.
After flattening, NESL can evenly distribute the computation workload
by splitting at places that user code doe not consider not sub-vector boundaries.

%%% \cite{DPH}
%%% \cite{Repa}
%%% \cite{Accelerate}

%%% \cite{Futhark}
%%% \cite{Nova}

The goal of isolating the cell-level portion of a program
has also led to some domain-specific languages,
targeting specific varieties of regular aggregate data.
Halide is a language designed for writing image processing pipelines \cite{Halide}.
The iteration space is the set of pixels in an image.
A Halide programmer writes code describing
how an individual pixel should be handled
(or a cluster of pixels for stencil computation),
and then in a separate portion of the program,
the programmer describes how the iteration space ought to be traversed.
Since the program does not intermingle
loop-nest control code with loop-body computation code,
it is easier for a maintainer to tune for performance
by adjusting the iteration schedule.

Diderot is a programming language
specifically aimed at processing medical images \cite{Diderot}.
The universe of data consists of tensor fields,
intended as functions on the continuous domain $\mathbb{R}^n$,
rather than any particular discrete index space.
Diderot offers pointwise arithmetic on tensors
and a collection of common operations such as outer product and transposition.
Aggregate lifting appears again,
albeit in a smaller form,
with some operations on tensor fields.

HorseIR uses implicit lifting over arrays
as an intermediate representation for SQL queries \cite{HorseIR}.
HorseIR is lower-level than APL itself,
serving as a three-address IR with vector instructions,
and the IR itself is designed to ease loop fusion.
In contrast with APL, all HorseIR arrays are vectors---%
there are no higher-rank arrays, and scalars are represented as unit-length vectors.
There is also a list datatype for handling heterogeneous aggregate data,
such as one row of a database table.

%% Since we've brought up databases, should we discuss KDB/Q?

%% typing arrays
Our type system's use of restricted dependent types
is inspired by Dependent ML \cite{DependentML}.
While Dependent ML is designed with the expectation
that the index language has a fully decidable theory,
Remora's index language does not \cite{Durnev}.
In subsequent work,
Dependent ML also focused on ensuring safety of array index accesses,
using singleton and range types to ensure
that numbers used for indexing fell within arrays' bounds \cite{DMLBounds}.

Others have applied established type system machinery
to an APL-like computation model.
Thatte's coercion semantics \cite{ThatteCoercion}
uses a form subtyping in which
scalar types are subtypes of aggregates,
and aggregate types in certain situations are subtypes of higher-dimensional aggregates.
The subtyping rules emit coercions which invoke functions such as {\tt map} and {\tt replicate},
automatically adapting scalar functions to aggregate data.
However, the restrictions on treating aggregate types as subtypes of larger aggregates
prohibit lifting for functions which expect non-scalar input
(\ie, those whose expected rank is greater than zero)
and lifting to unequal frame shapes
(\eg, vector-matrix addition).
Gibbons's embedding of in an extended Haskell \cite{APLicative}
%%%% "... to define transposition, which is needed in order to implement reranking"!?
extensively uses type-level programming
and defines much of the rank-polymorphic lifting machinery
in terms of transposition.

Single Assignment C (SAC) is a variant of C without mutation \cite{SAC}.
The primary iteration mechanism in SAC is the {\tt with} loop,
in which the programmer describes a traversal of the index space
and builds an output array an element at a time.
The lack of mutation pushes the programmer
to avoid writing loop-carried dependence,
much like in the rank-polymorphic model.
Translation of well-behaved APL programs to SAC
tends to be straightforward \cite{APLSAC},
though the resulting loop structure is
somewhat specific to the function being lifted.
A SAC variant called Qube introduces a Dependent ML-style type system
to ensure the safety of indexed array element accesses \cite{Qube}.
Qube retains SAC's {\tt with} loop-based programming model---%
relying on range types in contrast to
APL's and Remora's implicitly lifted function application semantics---%
and due to its C roots,
it does not support first-class functions
to the extent of permitting arrays to appear in function position.

%% Compilation
Since the syntactic structure of a line of APL code,
in particular the meaning of the juxtaposition of two terms,
depends on the meanings of names which appear in the code,
standard APL does not admit a fully static parsing algorithm.
Due to this and other idiosyncratic warts,
past efforts to compile APL have typically targetted
large but well-behaved subsets of the language.

A prominent exception is the APL\textbackslash3000 compiler,
which could produce machine code for individual statements
and used interpretation to manage inter-statement control flow \cite{APL3000Compiler}.
This allowed some intraprocedural optimization,
such as fold-unfold and map-map fusion%
\footnote{In the APL community,
  these transformations are respectively referred to as
``beating'' (imagine a metronome producing a stream of beats)
and ``dragging along'' (amassing a chain of
delayed scalar operations to apply in a single loop body).}
\cite{AbramsAPLMachine},
analogous to deforestation
which arose in mainstream functional languages \cite{BurstallDarlington, WadlerListless}.
% Internal note: p31 of Wadler's dissertation contrasts his work with Abrams
It also ensured that names' dynamic meanings would be available
by the time execution reached any statement which used them.

Weiss and Saal instead applied interprocedural data-flow analysis
to resolve the syntactic classes of variable names in APL code \cite{ParsingAPL}.
This analysis is not complete for APL itself
due to the possibility that reassignment of a variable name
will change how some line of code parses.
However, the authors found that
real-world code did not make use of
the full freedom to manipulate
the syntactic structure of an APL program
by dynamically reassigning variables,
suggesting that this is a misfeature which can be discarded for little cost.
However, the common line-at-a-time compilation style
used in other work
also served to ensure that
by the time execution reaches a line of code,
its variable names, and thus their syntactic classes, could be resolved.

Budd describes a compiler for an APL variant
in which the ambiguity in parsing is avoided by
declaring identifiers before use
and name resolution is simplified by adopting lexical scope.
\cite{BuddCompiler}.
%TODO: more about what Budd's compiler does and doesn't do

APEX parallelizes an APL dialect
with restrictions on many APL features deemed
incompatible with compilation
(such as producing a string representation of an arbitrary function)
or not strictly necessary for practical use (such as {\tt GOTO}) \cite{Apex}.
APEX's shape compatibility rules for implicit aggregate lifting
are less permissive than APL,
\eg, prohibiting most cases of vector-matrix addition.
Instead, both arguments must have the same rank
or at least one argument passed to a scalar-consuming function
be a scalar or singleton vector.

A more recent line of work has focused on
intermediate representations of array programs,
such as the Typed Array Intermediate Language,
which makes aggregate operations explicit
using an {\tt each} primitive operator \cite{APLTAIL}.
The type system tracks arrays' ranks,
which provides enough information to recognize
when an {\tt each} call is needed,
but it is not meant to ensure that
lifting a function application has a well-defined result
(\ie, shape incompatibility is still possible).
$\mathcal{L}_0$ plots out a loop fusion strategy
inspired by control flow graph reduction \cite{T2GraphReduction}.
The $\mathcal{L}_0$ compiler splits array programs
into kernels of fused aggregate operations
based on the applicability of $T_2$ graph reductions
(merging nodes $X$ and $Y$ if $X$ is $Y$'s sole predecessor)
to the program's data flow graph.
%TODO: extend this to cover Futhark (which L_0 eventually turned into)

%% Formal methods? (not really used in type checking, maybe mention in future work)
%%% \cite{Inez}
%%% \cite{Makanin, Karhumaki, Plandowski}

